\documentclass[final]{agujournal2019}
\usepackage{url}
\usepackage{lineno}
\usepackage{soul}

\usepackage{xcolor}
\definecolor{default-linkcolor}{HTML}{A50000}
\definecolor{default-filecolor}{HTML}{A50000}
\definecolor{default-citecolor}{HTML}{4077C0}
\definecolor{default-urlcolor}{HTML}{4077C0}

\usepackage{hyperref}
\hypersetup{
  hidelinks,
  breaklinks=true
}


\providecommand{\tightlist}{\setlength{\itemsep}{0pt}\setlength{\parskip}{0pt}}
\providecommand{\listfigurename}{List of Figures}
\providecommand{\listtablename}{List of Tables}
\newcounter{none}
\setcounter{none}{0}

\makeatletter
\@ifundefined{pandocbounded}{
  \usepackage{graphicx}
  \newcommand{\pandocbounded}[1]{\setbox0=\hbox{#1}\ifdim\wd0>\linewidth\resizebox{\linewidth}{!}{#1}\else#1\fi}
}{}
\makeatother

\renewcommand{\thefigure}{S\arabic{figure}}
\renewcommand{\thetable}{S\arabic{table}}
\renewcommand{\thesection}{S\arabic{section}}
\renewcommand{\theequation}{S\arabic{equation}}

\usepackage{longtable,booktabs,array}
\usepackage{calc}
\usepackage{etoolbox}
\makeatletter
\patchcmd\longtable{\par}{\if@noskipsec\mbox{}\fi\par}{}{}
\makeatother
\IfFileExists{footnotehyper.sty}{\usepackage{footnotehyper}}{\usepackage{footnote}}
\makesavenoteenv{longtable}

\NewDocumentCommand\citeproctext{}{}
\NewDocumentCommand\citeproc{mm}{\begingroup\def\citeproctext{#2}\cite{#1}\endgroup}
\makeatletter
 \let\@cite@ofmt\@firstofone
 \def\@biblabel#1{}
 \def\@cite#1#2{{#1\if@tempswa , #2\fi}}
\makeatother

\newlength{\cslhangindent}
\setlength{\cslhangindent}{1.5em}
\newlength{\csllabelwidth}
\setlength{\csllabelwidth}{3em}
\newenvironment{CSLReferences}[2]
 {\begin{list}{}{
  \setlength{\itemindent}{0pt}
  \setlength{\leftmargin}{0pt}
  \setlength{\parsep}{0pt}
  \ifodd #1
   \setlength{\leftmargin}{\cslhangindent}
   \setlength{\itemindent}{-\cslhangindent}
  \fi
  \setlength{\itemsep}{#2\baselineskip}
  }
  \item\relax}
 {\end{list}}
\usepackage{calc}
\newcommand{\CSLBlock}[1]{#1\hfill\break}
\newcommand{\CSLLeftMargin}[1]{\parbox[t]{\csllabelwidth}{#1}}
\newcommand{\CSLRightInline}[1]{\parbox[t]{\linewidth - \csllabelwidth}{#1}\break}
\newcommand{\CSLIndent}[1]{\hspace{\cslhangindent}#1}

\usepackage{amsmath}
\usepackage{amssymb}
\usepackage{amsfonts}
\makeatletter
\@ifpackageloaded{subcaption}{}{\usepackage{subcaption}}
\@ifpackageloaded{caption}{}{\usepackage{caption}}
\captionsetup[subfigure]{margin=0.5em}
\AtBeginDocument{%
\renewcommand*\figurename{Figure}
\renewcommand*\tablename{Table}
}
\AtBeginDocument{%
\renewcommand*\listfigurename{List of Figures}
\renewcommand*\listtablename{List of Tables}
}
\newcounter{pandoccrossref@subfigures@footnote@counter}
\newenvironment{pandoccrossrefsubfigures}{%
\setcounter{pandoccrossref@subfigures@footnote@counter}{0}
\begin{figure}\centering%
\gdef\global@pandoccrossref@subfigures@footnotes{}%
\DeclareRobustCommand{\footnote}[1]{\footnotemark%
\stepcounter{pandoccrossref@subfigures@footnote@counter}%
\ifx\global@pandoccrossref@subfigures@footnotes\empty%
\gdef\global@pandoccrossref@subfigures@footnotes{{##1}}%
\else%
\g@addto@macro\global@pandoccrossref@subfigures@footnotes{, {##1}}%
\fi}}%
{\end{figure}%
\addtocounter{footnote}{-\value{pandoccrossref@subfigures@footnote@counter}}
\@for\f:=\global@pandoccrossref@subfigures@footnotes\do{\stepcounter{footnote}\footnotetext{\f}}%
\gdef\global@pandoccrossref@subfigures@footnotes{}}
\@ifpackageloaded{float}{}{\usepackage{float}}
\floatstyle{ruled}
\@ifundefined{c@chapter}{\newfloat{codelisting}{h}{lop}}{\newfloat{codelisting}{h}{lop}[chapter]}
\floatname{codelisting}{Listing}
\newcommand*\listoflistings{\listof{codelisting}{List of Listings}}
\makeatother

\journalname{JGR: Solid Earth}

\makeatletter
\def\titlepageheadlinefont{\fontsize{9}{11}\selectfont}

\def\ps@headings{
  \def\@oddfoot{\centerline{\small --\the\c@page--}}
  \let\@evenfoot\@oddfoot
  \def\@oddhead{}
  \let\@evenhead\@oddhead
}
\ps@headings

\def\ps@front{
  \def\@oddhead{
    \vbox to 0pt{
      \vskip-30pt
      \centerline{\titlepageheadlinefont JGR: Solid Earth}
      \vss
    }
  }
  \let\@evenhead\@oddhead
  \def\@oddfoot{\centerline{\small --\the\c@page--}}
  \let\@evenfoot\@oddfoot
}

\let\oldtitle\title
\renewcommand{\title}[1]{\oldtitle{#1}\thispagestyle{front}}
\makeatother

\begin{document}


\title{Supplementary Information for "Beyond Equilibrium II: Ringwoodite
Decomposition Kinetics Amplify 660 km Discontinuity Depressions and
Offer an Independent Control on Slab Stagnation"}

\authors{Buchanan Kerswell\affil{1}, John Wheeler\affil{1}, Rene
Gassmöller\affil{2}, J. Huw Davies\affil{3}, Isabel
Papanagnou\affil{4}, Sanne Cottaar\affil{4}, David P. Dobson\affil{5}}

\affiliation{1}{Department of Earth, Oceans and Ecological Sciences,
University of Liverpool, Liverpool L69 3BX, UK}
\affiliation{2}{GEOMAR Helmholtz Centre for Ocean Research Kiel, Kiel,
Germany}
\affiliation{3}{School of Earth and Environmental Sciences, Cardiff
University, Park Place, Cardiff, CF10 3AT, UK}
\affiliation{4}{Bullard Laboratories, Department of Earth Sciences,
University of Cambridge, Cambridge CB3 0EZ, UK}
\affiliation{5}{Department of Earth Sciences, University College London,
London WC1E 6BS, UK}

\correspondingauthor{Buchanan Kerswell}{b.kerswell@liverpool.ac.uk}

\section*{Contents of this file}\label{sec:contents}
\addcontentsline{toc}{section}{Contents of this file}

\begin{enumerate}
\def\labelenumi{\arabic{enumi}.}
\tightlist
\item
  Figures \ref{fig:material-property-profiles} to
  \ref{fig:plume-composition-660-b50}
\item
  Table~\ref{tbl:depth-profile-summary}
\end{enumerate}

\clearpage

\section*{Introduction}\label{sec:introduction}
\addcontentsline{toc}{section}{Introduction}

This supplementary information provides five figures and one table that
support the material model, kinetic calibration, and quantitative
results described in the main text.

Figures \ref{fig:material-property-profiles} and
\ref{fig:thermodynamic-property-profiles} show the reference material
and thermodynamic properties used to initialize and drive the
compressible mantle flow simulations. These profiles give the density,
thermal expansivity, compressibility, and specific heat capacity of each
mantle phase across the transition zone, along with the Gibbs free
energy, entropy, and volume changes that govern the 410, 520, and 660
phase transitions.

Figure~\ref{fig:regression} presents the Monte Carlo regression used to
calibrate the ringwoodite decomposition kinetic rate equation against
the high-overpressure experimental dataset of Kubo et al. (2002). This
regression constrains the activation energy and growth-rate prefactor
within the uncertainties reported in the main text and justifies the
five-order-of-magnitude range of the kinetic prefactor \(Z_\mathrm{ri}\)
explored in our simulations.

Figures \ref{fig:slab-composition-660-b01} to
\ref{fig:plume-composition-660-b50} show slab and plume simulation
snapshots for scenarios with a moderate lower-mantle viscosity jump
(\(b\) = 50) compared to the no-jump case (\(b\) = 1) shown in the main
text. These snapshots show how the three kinetic regimes identified at
the 660 (quasi-equilibrium, intermediate, and stagnation) are affected
when an additional viscosity contrast is imposed at the base of the
transition zone (i.e., slabs and plumes tend to stagnate at the 660).

Table~\ref{tbl:depth-profile-summary} reports the kinetic prefactor
\(Z_\mathrm{ri}\), 660 displacement and width, maximum reaction rate,
and maximum descent or ascent velocity through the 660 for every plume
and slab simulation across all eight tested scenarios. These values
underlie the trends discussed in the main text.

\clearpage

\section*{Material Model (Figures S1--S3)}\label{sec:material-model}
\addcontentsline{toc}{section}{Material Model (Figures S1--S3)}

\clearpage

\begin{figure}
\centering
\includegraphics[width=1\linewidth,height=\textheight,keepaspectratio,alt={Reference material properties used in ASPECT simulations showing density (\textbackslash bar\{\textbackslash rho\}), thermal expansivity (\textbackslash bar\{\textbackslash alpha\}), compressibility (\textbackslash bar\{\textbackslash beta\}), and specific heat capacity (\textbackslash bar\{C\}\_p). Colored profiles are for pure Mg olivine (ol), wadsleyite (wd), ringwoodite (ri), and the post-spinel (ps) assemblage bridgmanite + periclase. Profiles are computed using the methods described in see Section 2.2.1 in the main text. Gray bands indicate the 410, 520, and 660 transitions.}]{../figs/adiabat/material-property-profiles.png}
\caption{Reference material properties used in ASPECT simulations
showing density (\(\bar{\rho}\)), thermal expansivity
(\(\bar{\alpha}\)), compressibility (\(\bar{\beta}\)), and specific heat
capacity (\(\bar{C}_p\)). Colored profiles are for pure Mg olivine (ol),
wadsleyite (wd), ringwoodite (ri), and the post-spinel (ps) assemblage
bridgmanite + periclase. Profiles are computed using the methods
described in see Section 2.2.1 in the main text. Gray bands indicate the
410, 520, and 660 transitions.}\label{fig:material-property-profiles}
\end{figure}

\begin{figure}
\centering
\includegraphics[width=0.8\linewidth,height=\textheight,keepaspectratio,alt={Reference thermodynamic properties used in our ASPECT simulations. Colored profiles are for pure Mg olivine (ol), wadsleyite (wd), ringwoodite (ri), and the post-spinel (ps) assemblage bridgmanite + periclase. Profiles are computed using the same methods as described in Figure~ (see Section 2.2.1 in the main text). Gray bands indicate the 410, 520, and 660 transitions.}]{../figs/adiabat/thermodynamic-property-profiles.png}
\caption{Reference thermodynamic properties used in our ASPECT
simulations. Colored profiles are for pure Mg olivine (ol), wadsleyite
(wd), ringwoodite (ri), and the post-spinel (ps) assemblage bridgmanite
+ periclase. Profiles are computed using the same methods as described
in Figure~\ref{fig:material-property-profiles} (see Section 2.2.1 in the
main text). Gray bands indicate the 410, 520, and 660
transitions.}\label{fig:thermodynamic-property-profiles}
\end{figure}

\begin{figure}
\centering
\includegraphics[width=0.68\linewidth,height=\textheight,keepaspectratio,alt={Monte Carlo regression of high-overpressure experimental data from Kubo et al. (2002). The best-fit activation energy E\^{}\{\textbackslash ast\} = 310 \textbackslash pm 233 kJ mol\^{}\{-1\} (2\textbackslash sigma) and growth-rate prefactor \textbackslash ln(K\^{}\{\textbackslash prime\}) = -29.8 \textbackslash pm 22.1 are both consistent with Kubo et al. (2002) within uncertainty. The large uncertainties in E\^{}\{\textbackslash ast\} and \textbackslash ln(K\^{}\{\textbackslash prime\}) reflect few data points with relatively high experimental T uncertainties. Our exploration over five orders of magnitude in Z\_\textbackslash mathrm\{ri\} roughly captures this broad uncertainty.}]{../figs/kubo/regression.png}
\caption{Monte Carlo regression of high-overpressure experimental data
from Kubo et al. (2002). The best-fit activation energy \(E^{\ast}\) =
310 \(\pm\) 233 kJ mol\(^{-1}\) (2\(\sigma\)) and growth-rate prefactor
\(\ln(K^{\prime})\) = -29.8 \(\pm\) 22.1 are both consistent with Kubo
et al. (2002) within uncertainty. The large uncertainties in
\(E^{\ast}\) and \(\ln(K^{\prime})\) reflect few data points with
relatively high experimental \(T\) uncertainties. Our exploration over
five orders of magnitude in \(Z_\mathrm{ri}\) roughly captures this
broad uncertainty.}\label{fig:regression}
\end{figure}

\clearpage

\section*{Simulation Snapshots: Slabs and Plumes (Figures
S4--S7)}\label{sec:simulation-snapshots}
\addcontentsline{toc}{section}{Simulation Snapshots: Slabs and Plumes
(Figures S4--S7)}

\clearpage

\begin{figure}
\centering
\includegraphics[width=1\linewidth,height=\textheight,keepaspectratio,alt={Reference slab simulations with moderate inflow velocity, intermediate slab strength, and no lower mantle viscosity jump (scenario 4; Table 1 in main text). Rows show stagnation (top row: Z\_\textbackslash mathrm\{ri\} = 1.6e-03 mol\^{}2 J\^{}\{-2\} s\^{}\{-1\}), intermediate (middle row: Z\_\textbackslash mathrm\{ri\} = 6.0e-01 mol\^{}2 J\^{}\{-2\} s\^{}\{-1\}), and quasi-equilibrium (bottom row: Z\_\textbackslash mathrm\{ri\} = 1.6e+03 mol\^{}2 J\^{}\{-2\} s\^{}\{-1\}) kinetic regimes after 100 Ma. Columns show dynamic temperature \textbackslash hat\{T\} (left column), log viscosity \textbackslash eta (middle column), and volume fraction of the post-spinel assemblage X\_\textbackslash mathrm\{ps\} (right column). Thick black lines (left column) highlight the 10\% and 90\% wadsleyite and post-spinel volume fraction contours used to define 660 displacement and width. The background grid pattern indicates ASPECT's adaptive mesh refinement near high thermal gradients and phase transitions. White trace (left column) indicates the representative depth profile used to extract data, and the white arrows indicate where the 660 structure is evaluated.}]{../figs/simulation/compositions/slab-set2-composition-b01-0100.png}
\caption{Reference slab simulations with moderate inflow velocity,
intermediate slab strength, and no lower mantle viscosity jump (scenario
4; Table 1 in main text). Rows show stagnation (top row:
\(Z_\mathrm{ri}\) = 1.6e-03 mol\(^2\) J\(^{-2}\) s\(^{-1}\)),
intermediate (middle row: \(Z_\mathrm{ri}\) = 6.0e-01 mol\(^2\)
J\(^{-2}\) s\(^{-1}\)), and quasi-equilibrium (bottom row:
\(Z_\mathrm{ri}\) = 1.6e+03 mol\(^2\) J\(^{-2}\) s\(^{-1}\)) kinetic
regimes after 100 Ma. Columns show dynamic temperature \(\hat{T}\) (left
column), log viscosity \(\eta\) (middle column), and volume fraction of
the post-spinel assemblage \(X_\mathrm{ps}\) (right column). Thick black
lines (left column) highlight the 10\% and 90\% wadsleyite and
post-spinel volume fraction contours used to define 660 displacement and
width. The background grid pattern indicates ASPECT's adaptive mesh
refinement near high thermal gradients and phase transitions. White
trace (left column) indicates the representative depth profile used to
extract data, and the white arrows indicate where the 660 structure is
evaluated.}\label{fig:slab-composition-660-b01}
\end{figure}

\clearpage

\begin{figure}
\centering
\includegraphics[width=1\linewidth,height=\textheight,keepaspectratio,alt={Slab simulations with moderate inflow velocity, intermediate slab strength, and a moderate lower mantle viscosity jump (scenario 0; Table 1 in main text). Description and annotations follow Figure~.}]{../figs/simulation/compositions/slab-set2-composition-b50-0100.png}
\caption{Slab simulations with moderate inflow velocity, intermediate
slab strength, and a moderate lower mantle viscosity jump (scenario 0;
Table 1 in main text). Description and annotations follow
Figure~\ref{fig:slab-composition-660-b01}.}\label{fig:slab-composition-660-b50}
\end{figure}

\clearpage

\begin{figure}
\centering
\includegraphics[width=1\linewidth,height=\textheight,keepaspectratio,alt={Reference plume simulations with moderate inflow velocity, intermediate plume strength, and no lower mantle viscosity jump (scenario 4; Table 1 in main text). Rows show ultra-sluggish (top row: Z\_\textbackslash mathrm\{ri\} = 1.6e-03 mol\^{}2 J\^{}\{-2\} s\^{}\{-1\}), intermediate (middle row: Z\_\textbackslash mathrm\{ri\} = 6.0e-01 mol\^{}2 J\^{}\{-2\} s\^{}\{-1\}), and fast (bottom row: Z\_\textbackslash mathrm\{ri\} = 1.6e+03 mol\^{}2 J\^{}\{-2\} s\^{}\{-1\}) kinetic conditions after 50 Ma. All other annotations follow Figure~.}]{../figs/simulation/compositions/plume-set2-composition-b01-0050.png}
\caption{Reference plume simulations with moderate inflow velocity,
intermediate plume strength, and no lower mantle viscosity jump
(scenario 4; Table 1 in main text). Rows show ultra-sluggish (top row:
\(Z_\mathrm{ri}\) = 1.6e-03 mol\(^2\) J\(^{-2}\) s\(^{-1}\)),
intermediate (middle row: \(Z_\mathrm{ri}\) = 6.0e-01 mol\(^2\)
J\(^{-2}\) s\(^{-1}\)), and fast (bottom row: \(Z_\mathrm{ri}\) =
1.6e+03 mol\(^2\) J\(^{-2}\) s\(^{-1}\)) kinetic conditions after 50 Ma.
All other annotations follow
Figure~\ref{fig:slab-composition-660-b01}.}\label{fig:plume-composition-660-b01}
\end{figure}

\clearpage

\begin{figure}
\centering
\includegraphics[width=1\linewidth,height=\textheight,keepaspectratio,alt={Plume simulations with moderate inflow velocity, intermediate plume strength, and a moderate lower mantle viscosity jump (scenario 0; Table 1 in main text). Description and annotations follow Figure~.}]{../figs/simulation/compositions/plume-set2-composition-b50-0050.png}
\caption{Plume simulations with moderate inflow velocity, intermediate
plume strength, and a moderate lower mantle viscosity jump (scenario 0;
Table 1 in main text). Description and annotations follow
Figure~\ref{fig:plume-composition-660-b01}.}\label{fig:plume-composition-660-b50}
\end{figure}

\clearpage

\section*{Structure of the 660: Displacement and
Width}\label{sec:structure-of-the-660}
\addcontentsline{toc}{section}{Structure of the 660: Displacement and
Width}

\clearpage

\begin{longtable}[]{@{}
  >{\raggedright\arraybackslash}p{(\linewidth - 14\tabcolsep) * \real{0.0583}}
  >{\raggedleft\arraybackslash}p{(\linewidth - 14\tabcolsep) * \real{0.0874}}
  >{\raggedleft\arraybackslash}p{(\linewidth - 14\tabcolsep) * \real{0.1553}}
  >{\raggedleft\arraybackslash}p{(\linewidth - 14\tabcolsep) * \real{0.1262}}
  >{\raggedleft\arraybackslash}p{(\linewidth - 14\tabcolsep) * \real{0.0583}}
  >{\raggedleft\arraybackslash}p{(\linewidth - 14\tabcolsep) * \real{0.2233}}
  >{\raggedleft\arraybackslash}p{(\linewidth - 14\tabcolsep) * \real{0.2233}}
  >{\raggedleft\arraybackslash}p{(\linewidth - 14\tabcolsep) * \real{0.0680}}@{}}
\caption{Summary of the kinetic prefactor \(Z_\mathrm{ri}\), 660
structure (displacement and width), maximum reaction rate
\(\dot{X}_\mathrm{max}\), and maximum vertical velocity through the 660
\(\vec{v}_\mathrm{max}\) evaluated in plume and slab simulations after
50 and 100 Ma of evolution, respectively. Kinetic regimes correspond to
1: quasi-equilibrium, 2: intermediate, 3: stagnation. Units are
\(Z_\mathrm{ri}\): mol\(^2\) J\(^{-2}\) s\(^{-1}\), displacement: km,
width: km, \(\dot{X}_\mathrm{max}\): Ma\(^{-1}\),
\(\vec{v}_\mathrm{max}\): cm yr\(^{-1}\). Note that displacement values
are signed: positive values indicate downward displacement of the 660
below the nominal equilibrium depth, while negative values indicate
upward displacement above the nominal equilibrium
depth.}\label{tbl:depth-profile-summary}\tabularnewline
\toprule\noalign{}
\begin{minipage}[b]{\linewidth}\raggedright
Type
\end{minipage} & \begin{minipage}[b]{\linewidth}\raggedleft
Scenario
\end{minipage} & \begin{minipage}[b]{\linewidth}\raggedleft
\(Z_\mathrm{ri}\)
\end{minipage} & \begin{minipage}[b]{\linewidth}\raggedleft
Displacement
\end{minipage} & \begin{minipage}[b]{\linewidth}\raggedleft
Width
\end{minipage} & \begin{minipage}[b]{\linewidth}\raggedleft
\(\dot{X}_\mathrm{max}\)
\end{minipage} & \begin{minipage}[b]{\linewidth}\raggedleft
\(\vec{v}_\mathrm{max}\)
\end{minipage} & \begin{minipage}[b]{\linewidth}\raggedleft
Regime
\end{minipage} \\
\midrule\noalign{}
\endfirsthead
\toprule\noalign{}
\begin{minipage}[b]{\linewidth}\raggedright
Type
\end{minipage} & \begin{minipage}[b]{\linewidth}\raggedleft
Scenario
\end{minipage} & \begin{minipage}[b]{\linewidth}\raggedleft
\(Z_\mathrm{ri}\)
\end{minipage} & \begin{minipage}[b]{\linewidth}\raggedleft
Displacement
\end{minipage} & \begin{minipage}[b]{\linewidth}\raggedleft
Width
\end{minipage} & \begin{minipage}[b]{\linewidth}\raggedleft
\(\dot{X}_\mathrm{max}\)
\end{minipage} & \begin{minipage}[b]{\linewidth}\raggedleft
\(\vec{v}_\mathrm{max}\)
\end{minipage} & \begin{minipage}[b]{\linewidth}\raggedleft
Regime
\end{minipage} \\
\midrule\noalign{}
\endhead
\bottomrule\noalign{}
\endlastfoot
plume & 0 & 1.6e+03 & -17.9 & 6.5 & 2.09e+06 & 2.41 & NA \\
plume & 0 & 2.2e+02 & -17.9 & 6.5 & 2.87e+05 & 2.41 & NA \\
plume & 0 & 3.1e+01 & -17.9 & 6.5 & 4.04e+04 & 2.41 & NA \\
plume & 0 & 4.3e+00 & -17.9 & 6.5 & 5.62e+03 & 2.41 & NA \\
plume & 0 & 6.0e-01 & -17.9 & 6.5 & 7.78e+02 & 2.40 & NA \\
plume & 0 & 8.3e-02 & -17.9 & 6.5 & 1.06e+02 & 2.38 & NA \\
plume & 0 & 1.2e-02 & -18.0 & 6.6 & 1.73e+01 & 2.38 & NA \\
plume & 0 & 1.6e-03 & -19.0 & 7.1 & 3.65e+00 & 2.47 & NA \\
plume & 1 & 1.6e+03 & -13.7 & 5.2 & 2.77e+05 & 2.92 & NA \\
plume & 1 & 2.2e+02 & -13.6 & 5.2 & 3.78e+04 & 2.92 & NA \\
plume & 1 & 3.1e+01 & -13.6 & 5.2 & 5.32e+03 & 2.91 & NA \\
plume & 1 & 4.3e+00 & -13.6 & 5.2 & 7.39e+02 & 2.92 & NA \\
plume & 1 & 6.0e-01 & -13.7 & 5.2 & 1.03e+02 & 2.91 & NA \\
plume & 1 & 8.3e-02 & -13.7 & 5.4 & 1.49e+01 & 2.87 & NA \\
plume & 1 & 1.2e-02 & -13.8 & 6.4 & 4.35e+00 & 3.18 & NA \\
plume & 1 & 1.6e-03 & -15.8 & 8.8 & 3.83e+00 & 3.60 & NA \\
plume & 2 & 1.6e+03 & -21.2 & 5.4 & 1.03e+06 & 3.71 & NA \\
plume & 2 & 2.2e+02 & -21.2 & 5.4 & 1.42e+05 & 3.71 & NA \\
plume & 2 & 3.1e+01 & -21.2 & 5.4 & 2.00e+04 & 3.71 & NA \\
plume & 2 & 4.3e+00 & -21.2 & 5.4 & 2.76e+03 & 3.72 & NA \\
plume & 2 & 6.0e-01 & -21.2 & 5.4 & 3.79e+02 & 3.70 & NA \\
plume & 2 & 8.3e-02 & -21.2 & 5.4 & 5.10e+01 & 3.67 & NA \\
plume & 2 & 1.2e-02 & -21.4 & 5.8 & 1.02e+01 & 3.78 & NA \\
plume & 2 & 1.6e-03 & -22.1 & 7.9 & 4.58e+00 & 4.01 & NA \\
plume & 3 & 1.6e+03 & -13.6 & 5.2 & 1.07e+04 & 4.81 & NA \\
plume & 3 & 2.2e+02 & -13.6 & 5.2 & 1.48e+03 & 4.81 & NA \\
plume & 3 & 3.1e+01 & -13.6 & 5.2 & 2.08e+02 & 4.81 & NA \\
plume & 3 & 4.3e+00 & -13.7 & 5.2 & 2.95e+01 & 4.81 & NA \\
plume & 3 & 6.0e-01 & -13.8 & 6.3 & 8.06e+00 & 5.14 & NA \\
plume & 3 & 8.3e-02 & -17.1 & 5.8 & 6.08e+01 & 4.65 & NA \\
plume & 3 & 1.2e-02 & -17.9 & 6.6 & 2.31e+01 & 4.91 & NA \\
plume & 3 & 1.6e-03 & -19.1 & 7.9 & 5.55e+00 & 5.22 & NA \\
plume & 4 & 1.6e+03 & -10.3 & 6.9 & 3.72e+05 & 4.46 & NA \\
plume & 4 & 2.2e+02 & -10.3 & 6.9 & 5.05e+04 & 4.44 & NA \\
plume & 4 & 3.1e+01 & -10.3 & 6.9 & 6.98e+03 & 4.41 & NA \\
plume & 4 & 4.3e+00 & -10.3 & 6.8 & 9.17e+02 & 4.30 & NA \\
plume & 4 & 6.0e-01 & -10.2 & 6.8 & 1.07e+02 & 4.06 & NA \\
plume & 4 & 8.3e-02 & -10.7 & 6.8 & 2.68e+01 & 4.04 & NA \\
plume & 4 & 1.2e-02 & -13.0 & 6.7 & 1.03e+01 & 3.71 & NA \\
plume & 4 & 1.6e-03 & -14.2 & 9.2 & 3.33e+00 & 3.12 & NA \\
plume & 5 & 1.6e+03 & -10.2 & 6.8 & 4.24e+05 & 4.96 & NA \\
plume & 5 & 2.2e+02 & -10.2 & 6.8 & 5.84e+04 & 4.95 & NA \\
plume & 5 & 3.1e+01 & -10.2 & 6.8 & 8.15e+03 & 4.79 & NA \\
plume & 5 & 4.3e+00 & -10.2 & 6.7 & 1.08e+03 & 4.81 & NA \\
plume & 5 & 6.0e-01 & -10.2 & 6.7 & 1.31e+02 & 4.58 & NA \\
plume & 5 & 8.3e-02 & -10.5 & 6.7 & 2.04e+01 & 4.24 & NA \\
plume & 5 & 1.2e-02 & -12.4 & 6.9 & 8.87e+00 & 3.95 & NA \\
plume & 5 & 1.6e-03 & -14.0 & 8.8 & 3.23e+00 & 3.42 & NA \\
plume & 6 & 1.6e+03 & -13.8 & 6.3 & 1.04e+06 & 7.16 & NA \\
plume & 6 & 2.2e+02 & -13.8 & 6.3 & 1.42e+05 & 7.15 & NA \\
plume & 6 & 3.1e+01 & -13.8 & 6.3 & 2.00e+04 & 7.13 & NA \\
plume & 6 & 4.3e+00 & -13.8 & 6.3 & 2.73e+03 & 7.01 & NA \\
plume & 6 & 6.0e-01 & -13.8 & 6.3 & 3.58e+02 & 6.49 & NA \\
plume & 6 & 8.3e-02 & -13.8 & 6.2 & 4.81e+01 & 6.21 & NA \\
plume & 6 & 1.2e-02 & -14.0 & 7.2 & 1.15e+01 & 6.18 & NA \\
plume & 6 & 1.6e-03 & -16.5 & 9.9 & 6.66e+00 & 5.39 & NA \\
plume & 7 & 1.6e+03 & -13.8 & 6.2 & 1.05e+06 & 8.34 & NA \\
plume & 7 & 2.2e+02 & -13.8 & 6.2 & 1.44e+05 & 8.35 & NA \\
plume & 7 & 3.1e+01 & -13.8 & 6.2 & 2.03e+04 & 8.29 & NA \\
plume & 7 & 4.3e+00 & -13.8 & 6.2 & 2.77e+03 & 8.17 & NA \\
plume & 7 & 6.0e-01 & -13.8 & 6.1 & 3.65e+02 & 7.57 & NA \\
plume & 7 & 8.3e-02 & -13.7 & 6.1 & 4.95e+01 & 7.48 & NA \\
plume & 7 & 1.2e-02 & -14.0 & 7.1 & 1.22e+01 & 7.28 & NA \\
plume & 7 & 1.6e-03 & -16.2 & 10.1 & 6.66e+00 & 6.41 & NA \\
slab & 0 & 1.6e+03 & 26.7 & 11.0 & 1.35e+02 & 0.49 & 3 \\
slab & 0 & 2.2e+02 & 27.3 & 11.5 & 1.75e+01 & 0.48 & 3 \\
slab & 0 & 3.1e+01 & 29.0 & 6.5 & 1.67e+00 & 0.32 & 3 \\
slab & 0 & 4.3e+00 & 31.1 & 6.1 & 6.85e-01 & 0.41 & 3 \\
slab & 0 & 6.0e-01 & 35.2 & 8.3 & 3.22e-01 & 0.30 & 3 \\
slab & 0 & 8.3e-02 & 40.7 & 11.8 & 1.58e-01 & 0.20 & 3 \\
slab & 0 & 1.2e-02 & 46.6 & 15.0 & 7.70e-02 & 0.14 & 3 \\
slab & 0 & 1.6e-03 & 51.4 & 16.9 & 4.12e-02 & 0.09 & 3 \\
slab & 1 & 1.6e+03 & 22.5 & 4.0 & 7.16e+00 & 0.61 & 1 \\
slab & 1 & 2.2e+02 & 27.1 & 11.3 & 1.66e+01 & 0.58 & 1 \\
slab & 1 & 3.1e+01 & 28.2 & 12.2 & 2.47e+00 & 0.56 & 1 \\
slab & 1 & 4.3e+00 & 30.1 & 7.1 & 1.24e+00 & 0.54 & 1 \\
slab & 1 & 6.0e-01 & 34.6 & 7.8 & 5.03e-01 & 0.46 & 1 \\
slab & 1 & 8.3e-02 & 42.9 & 13.6 & 2.24e-01 & 0.32 & 3 \\
slab & 1 & 1.2e-02 & 50.8 & 18.0 & 9.75e-02 & 0.20 & 3 \\
slab & 1 & 1.6e-03 & 58.7 & 21.4 & 5.00e-02 & 0.13 & 3 \\
slab & 2 & 1.6e+03 & 30.0 & 5.3 & 1.91e+01 & 0.94 & 1 \\
slab & 2 & 2.2e+02 & 30.4 & 4.6 & 3.15e+00 & 0.91 & 1 \\
slab & 2 & 3.1e+01 & 33.4 & 6.7 & 1.29e+00 & 0.87 & 1 \\
slab & 2 & 4.3e+00 & 39.7 & 9.4 & 9.06e-01 & 0.80 & 2 \\
slab & 2 & 6.0e-01 & 48.9 & 14.7 & 3.95e-01 & 0.61 & 2 \\
slab & 2 & 8.3e-02 & 60.9 & 22.0 & 1.52e-01 & 0.37 & 3 \\
slab & 2 & 1.2e-02 & 68.9 & 25.3 & 6.15e-02 & 0.20 & 3 \\
slab & 2 & 1.6e-03 & 74.9 & 24.7 & 3.81e-02 & 0.14 & 3 \\
slab & 3 & 1.6e+03 & 22.4 & 4.0 & 3.37e+01 & 1.07 & 1 \\
slab & 3 & 2.2e+02 & 22.5 & 4.0 & 6.08e+00 & 1.10 & 1 \\
slab & 3 & 3.1e+01 & 28.2 & 12.1 & 3.59e+00 & 1.05 & 1 \\
slab & 3 & 4.3e+00 & 30.4 & 6.8 & 2.00e+00 & 1.02 & 1 \\
slab & 3 & 6.0e-01 & 39.8 & 11.6 & 9.13e-01 & 0.99 & 2 \\
slab & 3 & 8.3e-02 & 52.9 & 19.1 & 3.38e-01 & 0.67 & 2 \\
slab & 3 & 1.2e-02 & 65.6 & 26.7 & 1.24e-01 & 0.36 & 2 \\
slab & 3 & 1.6e-03 & 80.5 & 33.7 & 5.13e-02 & 0.22 & 3 \\
slab & 4 & 1.6e+03 & 25.1 & 6.3 & 3.44e+00 & 1.55 & 1 \\
slab & 4 & 2.2e+02 & 28.2 & 6.0 & 2.14e+01 & 1.73 & 1 \\
slab & 4 & 3.1e+01 & 29.9 & 7.2 & 6.39e+00 & 1.87 & 1 \\
slab & 4 & 4.3e+00 & 33.1 & 7.0 & 2.24e+00 & 1.68 & 1 \\
slab & 4 & 6.0e-01 & 39.6 & 11.8 & 9.12e-01 & 1.11 & 2 \\
slab & 4 & 8.3e-02 & 44.0 & 13.4 & 2.17e-01 & 0.49 & 2 \\
slab & 4 & 1.2e-02 & 43.1 & 10.7 & 1.12e-01 & 0.31 & 3 \\
slab & 4 & 1.6e-03 & 48.4 & 10.8 & 6.82e-02 & 0.25 & 3 \\
slab & 5 & 1.6e+03 & 24.4 & 5.8 & 6.03e+01 & 2.05 & 1 \\
slab & 5 & 2.2e+02 & 24.7 & 6.0 & 7.22e+00 & 1.92 & 1 \\
slab & 5 & 3.1e+01 & 27.6 & 6.8 & 4.52e+00 & 1.73 & 1 \\
slab & 5 & 4.3e+00 & 30.0 & 7.1 & 2.22e+00 & 1.17 & 1 \\
slab & 5 & 6.0e-01 & 39.6 & 12.2 & 1.41e+00 & 1.71 & 2 \\
slab & 5 & 8.3e-02 & 52.4 & 20.0 & 6.38e-01 & 1.31 & 2 \\
slab & 5 & 1.2e-02 & 55.8 & 20.1 & 1.36e-01 & 0.46 & 2 \\
slab & 5 & 1.6e-03 & 57.8 & 16.4 & 7.96e-02 & 0.29 & 3 \\
slab & 6 & 1.6e+03 & 31.4 & 5.4 & 2.83e+01 & 3.21 & 1 \\
slab & 6 & 2.2e+02 & 32.5 & 6.2 & 7.40e+00 & 3.07 & 1 \\
slab & 6 & 3.1e+01 & 36.5 & 8.9 & 3.45e+00 & 2.85 & 1 \\
slab & 6 & 4.3e+00 & 44.0 & 12.2 & 1.84e+00 & 2.39 & 2 \\
slab & 6 & 6.0e-01 & 54.7 & 18.7 & 7.04e-01 & 1.44 & 2 \\
slab & 6 & 8.3e-02 & 55.7 & 17.8 & 1.36e-01 & 0.44 & 2 \\
slab & 6 & 1.2e-02 & 54.5 & 13.1 & 1.06e-01 & 0.36 & 3 \\
slab & 6 & 1.6e-03 & 68.4 & 13.4 & 7.39e-02 & 0.32 & 3 \\
slab & 7 & 1.6e+03 & 28.1 & 12.0 & 1.04e+02 & 2.29 & 1 \\
slab & 7 & 2.2e+02 & 28.4 & 12.3 & 1.38e+01 & 2.35 & 1 \\
slab & 7 & 3.1e+01 & 30.0 & 7.1 & 5.71e+00 & 2.19 & 1 \\
slab & 7 & 4.3e+00 & 34.7 & 7.9 & 2.25e+00 & 2.01 & 1 \\
slab & 7 & 6.0e-01 & 44.3 & 14.6 & 1.37e+00 & 2.09 & 2 \\
slab & 7 & 8.3e-02 & 60.3 & 23.8 & 6.14e-01 & 1.59 & 2 \\
slab & 7 & 1.2e-02 & 79.6 & 34.3 & 1.94e-01 & 0.91 & 2 \\
slab & 7 & 1.6e-03 & 78.3 & 26.1 & 7.63e-02 & 0.48 & 2 \\
\end{longtable}

\clearpage

\section*{References}\label{sec:references}
\addcontentsline{toc}{section}{References}

\protect\phantomsection\label{refs}
\begin{CSLReferences}{1}{0}
\bibitem[\citeproctext]{ref-kubo2002}
Kubo, T., Ohtani, E., Kato, T., Urakawa, S., Suzuki, A., Kanbe, Y., et
al. (2002). Mechanisms and kinetics of the post-spinel transformation in
{Mg\(_2\)SiO\(_4\)}. \emph{Physics of the Earth and Planetary
Interiors}, \emph{129}(1-2), 153--171.

\end{CSLReferences}


\end{document}
